\chapter{Amniotomy (Artificial Rupture of Membranes)}

\section*{Overview}
Amniotomy is the deliberate artificial rupture of the amniotic membranes performed during labor. It is used to augment labor, assess the amniotic fluid, and permit internal fetal monitoring.

\section*{Alternative Names}
Artificial rupture of membranes, ARM, Breaking the waters

\section*{Principle}
Using a sterile instrument such as an amnihook, the clinician creates a controlled opening in the amniotic membranes during a vaginal examination. Release of amniotic fluid may strengthen uterine contractions and accelerate labor. The procedure also allows inspection of the fluid and placement of internal fetal monitors.

\section*{History}
Amniotomy has been performed since the late 19th century and became a common intervention for labor augmentation during the 20th century.

\section*{Why It Is Done}
Amniotomy is ordered to induce or augment labor when progress is slow, to allow internal fetal heart rate or intrauterine pressure monitoring, and to examine the amniotic fluid for meconium. It may also shorten the first stage of labor.

\section*{Is This a Primary or Secondary Test or Procedure}
Secondary procedure used adjunctively during labor management rather than as a primary intervention.

\section*{How It Is Performed}
The clinician performs a sterile vaginal examination to confirm vertex presentation and intact membranes. A sterile instrument is used to rupture the membranes with a controlled release of fluid. Fetal heart rate is monitored before and after the procedure to detect cord compression.

\section*{Preparation}
Fetal presentation, gestational age, and the absence of active bleeding or infection are confirmed. Maternal blood type and group B streptococcus status should be known, and the fetal heart rate tracing is reviewed before the procedure.

\section*{Contraindications and Precautions}
Amniotomy is avoided when the presentation is not vertex, or with active genital herpes, placenta previa, or vasa previa. Relative contraindications include an unstable fetal lie and non-reassuring fetal status. Caution is needed when the presenting part is not well applied to the cervix because cord prolapse may follow.

\section*{Risks}
Umbilical cord prolapse, fetal heart rate decelerations, and ascending infection (chorioamnionitis) are the main risks. The amnihook may rarely cause a scalp injury or cephalohematoma.

\section*{Aftercare and Recovery}
Labor is closely monitored, and the duration of membrane rupture is tracked to guide the timing of delivery. The mother is observed for fever or other signs of infection, and group B streptococcus prophylaxis is given when indicated.

\section*{Results and Interpretation}
The outcome is assessed by labor progress, cervical dilation, and fetal status rather than by a laboratory value. Effective labor augmentation with a reassuring fetal heart rate tracing indicates success.

\section*{Expected Results}
Regular contractions with progressive cervical dilation and a reassuring fetal heart rate; amniotic fluid that is clear, without meconium.

\section*{Follow-up}
Continuous or intermittent fetal monitoring with reassessment of labor progress is continued, and delivery generally occurs within 12--24 hours of rupture. Cesarean delivery is performed if labor fails to progress or fetal compromise develops.

\section*{References}
\begin{enumerate}
  \item MedlinePlus. Amniotomy. U.S. National Library of Medicine; 2025.
  \item Current Medical Diagnosis and Treatment. McGraw-Hill; 2025.
\end{enumerate}

\clearpage
